\uuid{Vk4D}
\titre{Dérivée partielle via la règle de dérivation composée}
\niveau{L2}
\module{Analyse}
\chapitre{Fonctions de plusieurs variables}
\sousChapitre{Dérivation des fonctions composées}
\theme{Règle de dérivation composée, dérivée partielle}
\auteur{Maxime Nguyen}
\datecreate{2026-03-24}
\organisation{amscc}
\difficulte{2}

\contenu{
\texte{
  Soit $f$ une fonction de classe $\mathcal{C}^1$ sur $\mathbb{R}^2$ vérifiant $\forall t \in \mathbb{R},\; f(t, t^2) = 1$.
}
\begin{enumerate}
  \item \question{En utilisant la dérivée de la composée, montrer que $\dfrac{\partial f}{\partial x}(0,0) = 0$.}
    \reponse{
      On pose $\varphi(t) = f(t, t^2)$. Par la règle de dérivation composée :
      \[
        \varphi'(t) = \frac{\partial f}{\partial x}(t, t^2) + 2t\,\frac{\partial f}{\partial y}(t, t^2).
      \]
      Comme $\varphi(t) = 1$ pour tout $t$, on a $\varphi'(t) = 0$, donc :
      \[
        \frac{\partial f}{\partial x}(t, t^2) + 2t\,\frac{\partial f}{\partial y}(t, t^2) = 0.
      \]
      En particulier pour $t = 0$ :
      \[
        \frac{\partial f}{\partial x}(0,0) = 0.
      \]
    }
\end{enumerate}
}
